\documentclass[10pt,letterpaper]{article}
\usepackage[letterpaper,margin=0.72in,headheight=16pt,footskip=28pt]{geometry}
\usepackage[T1]{fontenc}
\usepackage{lmodern}
\usepackage{microtype}
\usepackage[dvipsnames,table]{xcolor}
\usepackage{graphicx}
\usepackage{booktabs}
\usepackage{tabularx}
\usepackage{longtable}
\usepackage{array}
\usepackage{ragged2e}
\usepackage{pdflscape}
\usepackage{enumitem}
\usepackage{parskip}
\usepackage{needspace}
\usepackage[most]{tcolorbox}
\usepackage{listings}
\usepackage{upquote}
\usepackage{textcomp}
\usepackage{fancyhdr}
\usepackage{titlesec}
\usepackage{lastpage}
\usepackage{xurl}
\usepackage[hidelinks]{hyperref}
\usepackage{bookmark}

\definecolor{Navy}{HTML}{123047}
\definecolor{Teal}{HTML}{0E6670}
\definecolor{LightTeal}{HTML}{EAF4F4}
\definecolor{Critical}{HTML}{A12828}
\definecolor{High}{HTML}{A45D00}
\definecolor{Standard}{HTML}{356A45}
\definecolor{Defense}{HTML}{5D6470}
\definecolor{Ink}{HTML}{20262D}
\definecolor{Muted}{HTML}{66717B}
\definecolor{Panel}{HTML}{F5F7F8}
\definecolor{RuleLine}{HTML}{D7DEE3}

\hypersetup{
  pdftitle={GitHub Advanced Security Supported-Language Secure Coding Standards Catalog},
  pdfauthor={Application Security Engineering},
  colorlinks=true,
  linkcolor=Navy,
  urlcolor=Teal,
  citecolor=Navy
}

\pagestyle{fancy}
\fancyhf{}
\lhead{\textcolor{Muted}{GHAS Secure Coding Standards}}
\rhead{\textcolor{Muted}{Enterprise Baseline v1.0}}
\cfoot{\textcolor{Muted}{\thepage\ / \pageref*{LastPage}}}
\renewcommand{\headrulewidth}{0.3pt}
\renewcommand{\footrulewidth}{0pt}

\titleformat{\section}{\Large\bfseries\color{Navy}}{\thesection}{0.65em}{}
\titleformat{\subsection}{\large\bfseries\color{Teal}}{\thesubsection}{0.6em}{}
\titleformat{\subsubsection}{\normalsize\bfseries\color{Ink}}{\thesubsubsection}{0.5em}{}
\titlespacing*{\section}{0pt}{1.7em}{0.7em}
\titlespacing*{\subsection}{0pt}{1.3em}{0.5em}
\setcounter{tocdepth}{2}
\setcounter{secnumdepth}{2}

\newcommand{\code}[1]{\texttt{\small #1}}
\newcommand{\prioritybadge}[2]{\tcbox[on line,boxrule=0pt,arc=1.2mm,boxsep=0.7mm,left=1mm,right=1mm,top=0.4mm,bottom=0.4mm,colback=#2!12,coltext=#2]{\scriptsize\bfseries #1}}
\newcommand{\pill}[1]{\tcbox[on line,boxrule=0pt,arc=1mm,boxsep=0.6mm,left=0.8mm,right=0.8mm,top=0.3mm,bottom=0.3mm,colback=LightTeal,coltext=Teal]{\scriptsize\bfseries #1}}

\newtcolorbox{rulebox}[1][]{
  enhanced, breakable, colback=white, colframe=RuleLine,
  boxrule=0.55pt, arc=1.4mm, left=2.5mm,right=2.5mm,top=1.5mm,bottom=1.5mm,
  borderline west={2.2pt}{0pt}{Teal},
  fonttitle=\bfseries\color{Navy}, coltitle=Navy,
  colbacktitle=Panel, attach boxed title to top left={xshift=1.5mm,yshift*=-1mm},
  boxed title style={boxrule=0pt,arc=1mm,colback=Panel},
  before skip=8pt, after skip=8pt, #1
}

\newtcolorbox{callout}[2][]{
  enhanced, breakable, colback=#2!5, colframe=#2!40,
  boxrule=0.6pt, arc=1.3mm, left=3mm,right=3mm,top=2mm,bottom=2mm,
  #1
}

\lstdefinestyle{securecode}{
  basicstyle=\ttfamily\footnotesize,
  columns=fullflexible,
  keepspaces=true,
  showstringspaces=false,
  breaklines=true,
  breakatwhitespace=false,
  frame=single,
  rulecolor=\color{RuleLine},
  backgroundcolor=\color{Panel},
  xleftmargin=2mm,
  xrightmargin=1mm,
  aboveskip=6pt,
  belowskip=8pt,
  tabsize=2
}
\lstset{style=securecode}

\newcolumntype{Y}{>{\RaggedRight\arraybackslash}X}
\newcolumntype{P}[1]{>{\RaggedRight\arraybackslash}p{#1}}
\renewcommand{\arraystretch}{1.15}
\setlength{\emergencystretch}{3em}


\begin{document}
\hypersetup{pdftitle={Enterprise Secure Workflow Standard --- GitHub Actions},pdfauthor={Application Security Engineering}}
\fancyhead[L]{\textcolor{Muted}{Enterprise Secure Workflow Standard --- GitHub Actions}}
\fancyhead[R]{\textcolor{Muted}{v1.0 --- 23 September 2026}}
\hypersetup{pageanchor=false}
\begin{titlepage}
\thispagestyle{empty}
\vspace*{0.85in}
{\color{Teal}\rule{1.2in}{3pt}}\\[0.32in]
{\Huge\bfseries\color{Navy}Enterprise Secure Workflow Standard --- GitHub Actions\par}
\vspace{0.18in}
{\Large\color{Teal}GitHub Advanced Security / CodeQL Workflow Enforcement Profile\par}
\vspace{0.55in}
\begin{callout}{Teal}
\textbf{Purpose.} Provide a standalone, traceable secure-coding standard that connects normative implementation rules to CWE, CodeQL detection, companion verification, and SSDLC enforcement.
\end{callout}
\vfill
\begin{tabularx}{\textwidth}{@{}P{1.45in}Y@{}}
\textbf{Standard version} & 1.0 \\
\textbf{Verified} & 23 September 2026 \\
\textbf{Document status} & Enterprise baseline \\
\textbf{Scope} & GitHub Actions workflows, reusable workflows, action metadata, runners, artifacts, caches, credentials, and CI/CD trust boundaries \\
\textbf{Source policy} & Official and authoritative primary sources only \\
\textbf{Normative terms} & MUST, MUST NOT, SHOULD, SHOULD NOT, MAY \\
\end{tabularx}
\vspace{0.55in}
\end{titlepage}
\hypersetup{pageanchor=true}
\pagenumbering{roman}
\tableofcontents
\clearpage
\pagenumbering{arabic}

\section{Executive Summary}
GitHub Actions is treated as a CI/CD execution and security domain, not as a programming language. This standard governs attacker-controlled repository metadata, pull-request code, workflow expressions, shell/script execution, \code{GITHUB\_TOKEN} permissions, secrets, environments, third-party actions, OIDC, runners, artifacts, caches, reusable workflows, and privileged/untrusted job boundaries.

\begin{callout}{Navy}
\textbf{Interpretation boundary.} A clean CodeQL Actions result does not prove a workflow is secure. Runner isolation, repository rulesets, environment protection, cloud trust policies, artifact provenance, and operational credential controls remain material.
\end{callout}

\section{Authority and Support Profile}
\begin{callout}{Teal}
\textbf{Classification:} Official Security Guidance\\
\textbf{Primary authority:} GitHub Secure Use Reference and GitHub CodeQL Actions query help\\
\textbf{CodeQL identifier:} \code{actions}\\
\textbf{Analyzed artifacts:} workflow and action-metadata YAML\\
\textbf{Verification date:} 23 September 2026
\end{callout}

\section{Workflow Security Domains}
\begin{itemize}[leftmargin=*]
\item Untrusted GitHub context, expression expansion, and \code{run:} script injection.
\item \code{pull\_request\_target} and other privileged trigger designs.
\item Minimum \code{GITHUB\_TOKEN} permissions and least privilege.
\item Secret minimization, environments, masking, rotation, and OIDC short-lived credentials.
\item Third-party action provenance and full-length commit-SHA pinning.
\item Reusable-workflow trust, self-hosted runner isolation, artifact/cache poisoning, and privileged/untrusted job boundaries.
\item Dependency review and provenance controls at the pull-request boundary.
\end{itemize}

\section{Normative Secure Workflow Requirements}
\begin{rulebox}[title={\texttt{SCS-ACTIONS-INJ-001} \hfill \prioritybadge{Critical Control}{Critical}}]
\textbf{Domain.} Injection Prevention

\textbf{Requirement.} Untrusted GitHub context values MUST NOT be interpolated directly into \code{run:} scripts; pass them through an intermediate environment variable and quote/use them as data for the target shell.

\textbf{Rationale.} Expression expansion happens before the shell executes and can turn pull-request metadata or other attacker-controlled fields into shell syntax.

\textbf{Risk.} Code / command injection

\begin{description}[leftmargin=2.65cm,labelwidth=2.3cm,style=nextline,itemsep=1pt,topsep=3pt]
  \item[CWE] CWE-94, CWE-95, CWE-116
  \item[OWASP] Context-dependent / not forced
  \item[CERT] Not mapped
  \item[CodeQL] \textbf{Direct}; actions/code-injection/critical
  \item[Enforcement] CodeQL required check + workflow review
  \item[Prevention] Structured/parameterized APIs; keep untrusted data out of interpreter syntax.
  \item[Verification] CodeQL required check + workflow review
  \item[Automation] Yes
  \item[Sources] \texttt{GH-SECURE-ACTIONS, GH-BUILTIN}
\end{description}
\end{rulebox}
\begin{rulebox}[title={\texttt{SCS-ACTIONS-AUTHZ-001} \hfill \prioritybadge{Critical Control}{Critical}}]
\textbf{Domain.} Authorization

\textbf{Requirement.} Each workflow MUST declare the minimum \code{GITHUB\_TOKEN} permissions required; write permissions MUST be granted only to jobs that need them.

\textbf{Rationale.} Over-privileged workflow tokens increase the blast radius of workflow or dependency compromise.

\textbf{Risk.} Excess privilege

\begin{description}[leftmargin=2.65cm,labelwidth=2.3cm,style=nextline,itemsep=1pt,topsep=3pt]
  \item[CWE] CWE-275, CWE-250
  \item[OWASP] Context-dependent / not forced
  \item[CERT] Not mapped
  \item[CodeQL] \textbf{Direct}; actions/missing-workflow-permissions
  \item[Enforcement] CodeQL + organization workflow policy
  \item[Prevention] Server-side, default-deny checks at protected object and operation boundaries.
  \item[Verification] CodeQL + organization workflow policy
  \item[Automation] Yes
  \item[Sources] \texttt{GH-SECURE-ACTIONS, GH-BUILTIN}
\end{description}
\end{rulebox}
\begin{rulebox}[title={\texttt{SCS-ACTIONS-DEP-001} \hfill \prioritybadge{Critical Control}{Critical}}]
\textbf{Domain.} Dependency / Package Security

\textbf{Requirement.} Third-party actions SHOULD be pinned to a full-length commit SHA; tag-only references require explicit risk acceptance and compensating review.

\textbf{Rationale.} Tags can move; a full commit SHA is immutable and reduces dependency substitution risk.

\textbf{Risk.} Untrusted dependency / action

\begin{description}[leftmargin=2.65cm,labelwidth=2.3cm,style=nextline,itemsep=1pt,topsep=3pt]
  \item[CWE] CWE-829
  \item[OWASP] Context-dependent / not forced
  \item[CERT] Not mapped
  \item[CodeQL] \textbf{Direct}; actions/unpinned-tag (security-extended/security-and-quality)
  \item[Enforcement] CodeQL security-extended + Actions policy + Dependabot
  \item[Prevention] Approved registries, reproducible dependency state, dependency review, and timely remediation.
  \item[Verification] CodeQL security-extended + Actions policy + Dependabot
  \item[Automation] Yes
  \item[Sources] \texttt{GH-SECURE-ACTIONS, GH-BUILTIN}
\end{description}
\end{rulebox}
\begin{rulebox}[title={\texttt{SCS-ACTIONS-SECRET-001} \hfill \prioritybadge{Critical Control}{Critical}}]
\textbf{Domain.} Secrets

\textbf{Requirement.} Workflow secrets MUST be minimized, scoped to the smallest job/environment, never echoed, and MUST NOT be exposed to untrusted fork/pull-request code; cloud authentication SHOULD use OIDC short-lived credentials where supported.

\textbf{Rationale.} Workflow logs, subprocesses, third-party actions, and untrusted code can exfiltrate long-lived secrets.

\textbf{Risk.} Secret exposure / credential theft

\begin{description}[leftmargin=2.65cm,labelwidth=2.3cm,style=nextline,itemsep=1pt,topsep=3pt]
  \item[CWE] CWE-312, CWE-532
  \item[OWASP] Context-dependent / not forced
  \item[CERT] Not mapped
  \item[CodeQL] \textbf{Direct/Partial}; Actions excessive-secret-exposure and unmasked-secret queries
  \item[Enforcement] CodeQL + secret scanning + environments + OIDC
  \item[Prevention] External secret stores, least privilege, and rotation/revocation after exposure.
  \item[Verification] CodeQL + secret scanning + environments + OIDC
  \item[Automation] Partial
  \item[Sources] \texttt{GH-SECURE-ACTIONS, GH-OIDC, GH-SECRET, GH-BUILTIN}
\end{description}
\end{rulebox}
\begin{rulebox}[title={\texttt{SCS-ACTIONS-AUTHZ-002} \hfill \prioritybadge{Critical Control}{Critical}}]
\textbf{Domain.} Authorization

\textbf{Requirement.} \code{pull\_request\_target} or other privileged triggers MUST NOT check out and execute untrusted pull-request code with write tokens or secrets.

\textbf{Rationale.} Privileged base-repository context combined with attacker-controlled code creates a repository compromise path.

\textbf{Risk.} Untrusted code with elevated privileges

\begin{description}[leftmargin=2.65cm,labelwidth=2.3cm,style=nextline,itemsep=1pt,topsep=3pt]
  \item[CWE] CWE-829, CWE-269
  \item[OWASP] Context-dependent / not forced
  \item[CERT] Not mapped
  \item[CodeQL] \textbf{Direct/Partial}; Actions untrusted-checkout/privileged-trigger queries
  \item[Enforcement] CodeQL + workflow architecture review
  \item[Prevention] Server-side, default-deny checks at protected object and operation boundaries.
  \item[Verification] CodeQL + workflow architecture review
  \item[Automation] Partial
  \item[Sources] \texttt{GH-SECURE-ACTIONS, GH-BUILTIN}
\end{description}
\end{rulebox}
\begin{rulebox}[title={\texttt{SCS-ACTIONS-BUILD-001} \hfill \prioritybadge{High Priority}{High}}]
\textbf{Domain.} Compiler / Build Security

\textbf{Requirement.} Security workflows SHOULD use CodeQL default/code-scanning queries plus \code{security-extended} where the organization accepts the additional signal, and required checks SHOULD block merges according to documented severity policy.

\textbf{Rationale.} Some security queries, including unpinned-action detection and representative Rust lifetime analysis, are not in the default code-scanning suite.

\textbf{Risk.} Coverage gap from suite selection

\begin{description}[leftmargin=2.65cm,labelwidth=2.3cm,style=nextline,itemsep=1pt,topsep=3pt]
  \item[CWE] Not forced / context-dependent
  \item[OWASP] Context-dependent / not forced
  \item[CERT] Not mapped
  \item[CodeQL] \textbf{Requires configuration}; Query-suite configuration
  \item[Enforcement] CodeQL workflow + branch/ruleset required checks
  \item[Prevention] Supported toolchains, reproducible builds, static analysis, and security diagnostics in CI.
  \item[Verification] CodeQL workflow + branch/ruleset required checks
  \item[Automation] Yes
  \item[Sources] \texttt{GH-BUILTIN, GH-WORKFLOW}
\end{description}
\end{rulebox}
\begin{rulebox}[title={\texttt{SCS-ACTIONS-FILE-001} \hfill \prioritybadge{High Priority}{High}}]
\textbf{Domain.} File / Path Security

\textbf{Requirement.} Artifacts and caches from untrusted or lower-trust jobs MUST NOT be consumed by privileged jobs as executable/configuration input without integrity and provenance controls.

\textbf{Rationale.} Artifact/cache poisoning can move attacker-controlled content across job or privilege boundaries.

\textbf{Risk.} Artifact/cache poisoning

\begin{description}[leftmargin=2.65cm,labelwidth=2.3cm,style=nextline,itemsep=1pt,topsep=3pt]
  \item[CWE] CWE-829, CWE-94
  \item[OWASP] Context-dependent / not forced
  \item[CERT] Not mapped
  \item[CodeQL] \textbf{Direct/Partial}; Actions artifact-poisoning / cache-poisoning query families
  \item[Enforcement] CodeQL + workflow isolation + provenance checks
  \item[Prevention] Root confinement, safe temporary/archive APIs, and handle-based authorization where possible.
  \item[Verification] CodeQL + workflow isolation + provenance checks
  \item[Automation] Partial
  \item[Sources] \texttt{GH-SECURE-ACTIONS, GH-BUILTIN}
\end{description}
\end{rulebox}
\begin{rulebox}[title={\texttt{SCS-ACTIONS-DEP-002} \hfill \prioritybadge{High Priority}{High}}]
\textbf{Domain.} Dependency / Package Security

\textbf{Requirement.} Pull requests that change dependency manifests/lockfiles or action references SHOULD run dependency review and MUST block additions that violate the organization's vulnerability/license policy.

\textbf{Rationale.} Dependency review exposes dependency changes at the pull-request boundary, before merge.

\textbf{Risk.} Vulnerable dependency introduction

\begin{description}[leftmargin=2.65cm,labelwidth=2.3cm,style=nextline,itemsep=1pt,topsep=3pt]
  \item[CWE] CWE-829
  \item[OWASP] Context-dependent / not forced
  \item[CERT] Not mapped
  \item[CodeQL] \textbf{Requires another control}; Dependency Review Action
  \item[Enforcement] Dependency review required check + Dependabot
  \item[Prevention] Approved registries, reproducible dependency state, dependency review, and timely remediation.
  \item[Verification] Dependency review required check + Dependabot
  \item[Automation] Yes
  \item[Sources] \texttt{GH-DEP-REVIEW, GH-DEPENDABOT}
\end{description}
\end{rulebox}
\section{Verified GitHub Actions CodeQL Crosswalk}\begin{landscape}\scriptsize
\setlength{\tabcolsep}{2.6pt}
\begin{longtable}{@{}P{1.28in}P{1.12in}P{1.52in}P{0.78in}P{0.52in}P{0.48in}P{1.32in}P{0.82in}@{}}
\toprule
\textbf{Rule ID(s)} & \textbf{Weakness} & \textbf{Query ID} & \textbf{CWE} & \textbf{Prec.} & \textbf{Sec.} & \textbf{Suite(s)} & \textbf{Coverage}\\
\midrule
\endfirsthead
\toprule
\textbf{Rule ID(s)} & \textbf{Weakness} & \textbf{Query ID} & \textbf{CWE} & \textbf{Prec.} & \textbf{Sec.} & \textbf{Suite(s)} & \textbf{Coverage}\\
\midrule
\endhead
SCS-ACTIONS-INJ-001 & Code injection & \href{https://codeql.github.com/codeql-query-help/actions/actions-code-injection-critical/}{\code{actions/code-injection/critical}} & CWE-94, CWE-95, CWE-116 & very-high & 9.0 & code-scanning; sec.-extended; sec.-and-quality & Direct example\\
SCS-ACTIONS-AUTHZ-001 & Missing workflow permissions & \href{https://codeql.github.com/codeql-query-help/actions/actions-missing-workflow-permissions/}{\code{actions/missing-workflow-permissions}} & CWE-275 & high & 5.0 & code-scanning; sec.-extended; sec.-and-quality & Direct example\\
SCS-ACTIONS-DEP-001 & Unpinned action tag & \href{https://codeql.github.com/codeql-query-help/actions/actions-unpinned-tag/}{\code{actions/unpinned-tag}} & CWE-829 & medium & 5.0 & sec.-extended; sec.-and-quality & Direct example\\
\bottomrule
\end{longtable}
\normalsize\end{landscape}
\section{Companion Enforcement}
\begin{tabularx}{\textwidth}{@{}P{1.45in}Y P{2.3in}@{}}
\toprule
\textbf{Area} & \textbf{Primary prevention} & \textbf{Companion verification / governance}\\
\midrule
Script injection & Intermediate environment variables; structured action inputs; quote shell data & Actions CodeQL + workflow review\\
Permissions & Explicit least-privilege \code{permissions} & Repository/org policy + ruleset review\\
Secrets & Minimum-scope environments/secrets; OIDC where supported & Secret scanning + rotation + cloud trust-policy review\\
Third-party actions & Full-length immutable commit SHA & Dependabot/dependency review + provenance review\\
Privileged triggers & Separate trusted and untrusted execution; never execute PR code in privileged context & Architecture/workflow review + tests\\
Runners & Ephemeral/isolated runners for untrusted code where appropriate & Runner hardening and lifecycle governance\\
Artifacts/caches & Treat lower-trust outputs as untrusted; validate provenance/integrity before privileged consumption & CodeQL query families + workflow isolation tests\\
\bottomrule
\end{tabularx}

\section{SSDLC Integration}
\small
\begin{tabularx}{\textwidth}{@{}P{1.12in}Y P{2.15in}@{}}
\toprule
\textbf{Stage} & \textbf{Engineering activity} & \textbf{Required evidence / gate}\\
\midrule
Design & Identify trust boundaries, sensitive data, authorization model, external dependencies, native/unsafe boundaries, and abuse cases. & Threat/architecture review where risk warrants.\\
Implementation & Apply applicable SCS-ALL and language overlay rules. & Code review; compiler/type/lint diagnostics.\\
Pull Request & Review code, dependency, workflow, and secret changes. & Dependency review; push protection; reviewer approval.\\
CodeQL & Run the configured language queries and required suites. & No unaccepted merge-blocking findings.\\
Build / Test & Run language-native security checks and negative-path tests. & CI status checks, race/sanitizer/fuzz results as applicable.\\
Security Testing & Exercise runtime behavior not reliably proven statically. & DAST/IAST/integration/property/fuzz testing as appropriate.\\
Merge Gate & Enforce repository/ruleset policy. & Required checks + approved exceptions only.\\
Deployment & Preserve secure configuration, credentials, TLS, and least privilege. & Deployment policy/configuration evidence.\\
Runtime Monitoring & Capture security-relevant events without secrets. & Central monitoring and alerting.\\
Remediation & Fix, validate, merge, and close findings with traceability. & Retest/rescan; time-bounded risk acceptance if unresolved.\\
\bottomrule
\end{tabularx}
\normalsize

\section{Workflow Pre-Merge Checklist}
\begin{itemize}[leftmargin=*,itemsep=2pt]
\item Untrusted context values are not interpolated directly into executable \code{run:} scripts.
\item \code{GITHUB\_TOKEN} permissions are explicitly minimized.
\item Privileged triggers do not check out/execute untrusted pull-request code with secrets or write tokens.
\item Third-party actions are pinned to approved immutable revisions, preferably full-length commit SHAs.
\item Secrets are minimized, scoped to environments/jobs, never echoed, and replaced with OIDC short-lived credentials where supported.
\item Self-hosted runners, reusable workflows, caches, and artifacts have explicit trust-boundary controls.
\item Dependency review and required CodeQL/workflow checks pass before merge.
\end{itemize}

\begin{landscape}
\section{Official Reference Catalog}
All references below are first-party or authoritative sources used by this standalone standard. Access / verification date: 23 September 2026.

\scriptsize
\setlength{\tabcolsep}{3pt}
\begin{longtable}{@{}P{0.82in}P{1.23in}P{2.25in}P{1.7in}P{3.0in}@{}}
\toprule
\textbf{Key} & \textbf{Organization} & \textbf{Document} & \textbf{Applies to} & \textbf{URL}\\
\midrule
\endfirsthead
\toprule
\textbf{Key} & \textbf{Organization} & \textbf{Document} & \textbf{Applies to} & \textbf{URL}\\
\midrule
\endhead
GH-BUILTIN & GitHub & CodeQL built-in queries & Built-in CodeQL query inventories & \url{https://docs.github.com/en/code-security/reference/code-scanning/codeql/codeql-queries/built-in-queries}\\
GH-DEP-REVIEW & GitHub & About dependency review & Dependency review and pull-request gating & \url{https://docs.github.com/en/code-security/supply-chain-security/understanding-your-software-supply-chain/about-dependency-review}\\
GH-DEPENDABOT & GitHub & About Dependabot security updates & Vulnerable dependency remediation & \url{https://docs.github.com/en/code-security/dependabot/dependabot-security-updates/about-dependabot-security-updates}\\
GH-LANG & GitHub & CodeQL --- Supported languages and frameworks & All CodeQL languages; versions, compilers, frameworks, libraries & \url{https://codeql.github.com/docs/codeql-overview/supported-languages-and-frameworks/}\\
GH-OIDC & GitHub & About security hardening with OpenID Connect & Short-lived cloud credentials from Actions & \url{https://docs.github.com/en/actions/concepts/security/openid-connect}\\
GH-PUSH & GitHub & About push protection & Blocking supported secrets before repository exposure & \url{https://docs.github.com/en/code-security/secret-scanning/introduction/about-push-protection}\\
GH-SECRET & GitHub & About secret scanning & Secret discovery & \url{https://docs.github.com/en/code-security/secret-scanning/introduction/about-secret-scanning}\\
GH-SECURE-ACTIONS & GitHub & Secure use reference --- GitHub Actions & GitHub Actions workflow security & \url{https://docs.github.com/en/actions/reference/security/secure-use}\\
GH-WORKFLOW & GitHub & Configuring CodeQL workflow / language identifiers and build modes & CodeQL language identifiers and compiled-language build behavior & \url{https://docs.github.com/en/code-security/code-scanning/creating-an-advanced-setup-for-code-scanning/configuring-codeql-workflow-for-compiled-languages}\\
MITRE-CWE & MITRE & Common Weakness Enumeration (CWE) & Weakness taxonomy and mappings & \url{https://cwe.mitre.org/}\\
NIST-SSDF & NIST & SP 800-218, Secure Software Development Framework (SSDF) Version 1.1 & Cross-language secure software development practices & \url{https://csrc.nist.gov/pubs/sp/800/218/final}\\
OWASP-ASVS & OWASP & Application Security Verification Standard (ASVS) 5.0 & Web/application security verification requirements & \url{https://owasp.org/www-project-application-security-verification-standard/}\\
\bottomrule
\end{longtable}
\normalsize
\end{landscape}

\appendix
\section{Build and Maintenance}
\begin{lstlisting}
latexmk -lualatex -interaction=nonstopmode -halt-on-error secure-workflow-standard-github-actions.tex
\end{lstlisting}
Revalidate Actions query metadata, runner guidance, token semantics, and supply-chain guidance whenever GitHub changes the Actions security model or CodeQL bundle.
\end{document}
